\clearpage
\subsection*{Can the trained models answer it?}
We locked five predictions in git before training on fresh seeds, then scored them with
a script, first for the original circuit and then for the improved v2 circuit
(eight copies of the same design; see the next section for how it was built).

\begin{center}
\small
\begin{tabular}{l c c}
\toprule
Prediction & original circuit & v2 circuit \\
\midrule
G-1\quad error below 1\% & no (53\%) & \textbf{yes (0.09\%)} \\
G-2\quad beats a same-size classical network & no & \textbf{yes (18$\times$)} \\
G-3\quad beats the same circuit with random frequencies & no & \textbf{yes (9.9$\times$)} \\
G-4\quad improvement lies on the designed frequencies & \textbf{yes (\GWVOneOverlap\%)} & \textbf{yes (99.4\%)} \\
G-5\quad waterlogged length within 5\% of the truth & no (0\,m) & \textbf{yes (714\,m vs 713\,m)} \\
\midrule
Held & \GWVOneConf{} of 5 & \GWVTwoConf{} of 5 \\
\bottomrule
\end{tabular}
\end{center}

\begin{center}
  \includegraphics[width=0.85\linewidth]{groundwater_models_v2}
\end{center}
The v2 circuit tracks the true water table to within a few millimetres across the whole
field and gets the waterlogged length right to 1\,m. The original circuit (below) missed
it entirely. The classical model handed the same frequencies fails here too.
\begin{center}
  \includegraphics[width=0.7\linewidth]{groundwater_models_v1}
\end{center}
